\section{FPGA 原型适配与实现结果}

\subsection{原型目标与实现路径}

为说明设计具备工程落地能力，本文进一步给出 FPGA 原型适配路径。
当前工作重点是先完成综合、实现、时序、约束和 bring-up 路径整理，
在此基础上确认该方案能够稳定映射到目标原型平台，再进入实体板验证。
本设计以 Nexys A7-100T 为目标平台，并以 Vivado \texttt{impl50} 流程作为统一实现条件。

\begin{figure}[H]
    \centering
    \begin{tikzpicture}[node distance=9mm and 9mm, every node/.style={font=\small}]
        \node[reportbox, minimum width=27mm] (rtl) {RTL / 脚本\\\scriptsize 设计输入};
        \node[reportbox, minimum width=27mm, right=of rtl] (impl) {Vivado impl50\\\scriptsize 50MHz 原型实现};
        \node[reportbox, minimum width=27mm, right=of impl] (rpt) {报告输出\\\scriptsize 资源 / 时序 / bit};
        \node[reportbox, minimum width=27mm, right=of rpt] (probe) {FPGA-like probe\\\scriptsize CoreMark 近板级观察};
        \node[reportbox, fill=orange!10, minimum width=36mm, below=16mm of rpt] (closure) {原型适配阶段\\\scriptsize 适配结果};
        \node[reportbox, fill=red!6, minimum width=36mm, below=16mm of probe] (board) {后续板级验证\\\scriptsize 实体板卡 / I/O 约束};

        \draw[reportline] (rtl) -- (impl);
        \draw[reportline] (impl) -- (rpt);
        \draw[reportline] (rpt) -- (probe);
        \draw[reportline] (impl.south) -- (closure.north west);
        \draw[reportline] (rpt.south) -- (closure.north east);
        \draw[reportdash] (probe.south) -- (board.north);
    \end{tikzpicture}
    \caption{FPGA 原型适配路径与板级边界示意图}
\end{figure}


\subsection{实现内容}

当前已经完成的 FPGA 原型适配工作包括：
\begin{itemize}
    \item Vivado \texttt{impl50} 构建脚本、bitstream 路径和报告目录整理完毕；
    \item FPGA 顶层、XDC、bring-up checklist 和相关 README 已统一到同一实现配置；
    \item \texttt{IMEM\_OUTPUT\_REG = 0} 与 \texttt{DMEM\_OUTPUT\_REG = 0} 的原型路径配置已固定；
    \item FPGA-like probe 已用于辅助观察原型路径下的运行周期与相对效率。
\end{itemize}

\subsection{关键结果}

\begin{table}[H]
    \centering
    \small
    \caption{FPGA 原型适配结果}
    \begin{tabularx}{0.95\textwidth}{L{0.25\textwidth} L{0.22\textwidth} Y}
        \toprule
        项目 & 结果 & 说明 \\
        \midrule
        目标平台 & Nexys A7-100T & 原型验证目标板卡 \\
        实现条件 & \texttt{impl50} & 50MHz 目标频率下的综合实现流程 \\
        Slice LUTs & \texttt{2556} & 综合实现后的逻辑资源占用 \\
        Slice Registers & \texttt{2170} & 综合实现后的寄存器资源占用 \\
        Block RAM Tile & \texttt{4} & 片上存储资源占用 \\
        DSP & \texttt{0} & 未依赖 DSP 资源 \\
        WNS & \texttt{+5.822ns} & 50MHz 目标下保持正建立时间裕量 \\
        WHS & \texttt{+0.057ns} & 保持正保持时间裕量 \\
        原型观测周期 & \texttt{156442 cycles} & 来自 FPGA-like probe 的近板级路径观测 \\
        原型观测 Score (CoreMark/MHz) & \texttt{7.728811} & 用于辅助观察原型路径下的相对效率 \\
        \bottomrule
    \end{tabularx}
\end{table}

\subsection{实体板相关工作}

这里给出的 FPGA 结果属于原型适配阶段结论，主要用于说明设计具备原型实现能力。
尚未完成的内容包括 UART/LED 实体连线验证、boot banner 现场观测以及最终板级 I/O 约束确认。
这些工作依赖实体板卡与现场条件，因此被保留为后续阶段任务。

\subsection{小结}

从现有结果看，\texttt{YH\_rv\_cpu} 已经完成 FPGA 原型适配的关键前置工作，
具备继续进入实体板 bring-up 的工程基础。对于初赛阶段而言，这说明该方案不仅完成了 RTL 设计，
也具备明确的原型实现路径。
